%=====================================================

% lista de acrônimos (siglas e abreviações)

\begin{listaabrev}

\begin{longtable}[l]{p{0.22\linewidth}p{0.70\linewidth}}
ABNT & Associação Brasileira de Normas Técnicas\\
ANN & \textit{Artificial Neural Network} (rede neural artificial)\\
CI & Circuito integrado\\
CISPR & \textit{Comité International Spécial des Perturbations Radioélectriques}\\
DE & \textit{Differential Evolution} (evolução diferencial)\\
EDA & \textit{Electronic Design Automation}\\
EMC & \textit{Electromagnetic Compatibility} (compatibilidade eletromagnética)\\
EMI & \textit{Electromagnetic Interference} (interferência eletromagnética)\\
GBM & \textit{Gradient Boosting Machine}\\
GPR & \textit{Gaussian Process Regression} (regressão por processos gaussianos)\\
IEC & \textit{International Electrotechnical Commission}\\
IEEE & \textit{Institute of Electrical and Electronics Engineers}\\
LHS & \textit{Latin Hypercube Sampling} (amostragem por hipercubo latino)\\
MAPE & \textit{Mean Absolute Percentage Error}\\
ML & \textit{Machine Learning} (aprendizado de máquina)\\
MLP & \textit{Multilayer Perceptron} (perceptron multicamadas)\\
MOEA & \textit{Multi-Objective Evolutionary Algorithm}\\
NSGA-II & \textit{Non-dominated Sorting Genetic Algorithm II}\\
PB & \textit{Physics-Based} (baseado em física)\\
PCB & \textit{Printed Circuit Board} (placa de circuito impresso)\\
PDN & \textit{Power Delivery Network} (rede de distribuição de energia)\\
PI & \textit{Power Integrity} (integridade de potência)\\
PINN & \textit{Physics-Informed Neural Network}\\
PPGEE & Programa de Pós-Graduação em Engenharia Elétrica\\
RF & \textit{Random Forest} (floresta aleatória)\\
RMSE & \textit{Root Mean Square Error} (raiz do erro quadrático médio)\\
SHAP & \textit{SHapley Additive exPlanations}\\
SI & \textit{Signal Integrity} (integridade de sinal)\\
SDD & \textit{Spec-Driven Development} (desenvolvimento orientado a especificação)\\
TCC & Trabalho de Conclusão de Curso\\
TDD & \textit{Test-Driven Development} (desenvolvimento orientado a testes)\\
TUHH & \textit{Technische Universität Hamburg}\\
UFPR & Universidade Federal do Paraná\\
VNA & \textit{Vector Network Analyzer} (analisador vetorial de redes)\\
\end{longtable}

\end{listaabrev}

%=====================================================
